\section{Spielvarianten}
Ich habe im Rahmen meines Exposés bereits eine grobe Schätzung darüber abgegeben, welche Pokervarianten ich lösen zu können glaube.
Dies beinhaltet zuerst das wohl einfachste Pokerspiel, Kuhn Poker.
Anschließend die häufig genutze Form Leduc Hold'em, welche genau wie Kuhn Poker, ausschließlich zu Forschungszwecken verwendet wird. 
Das dritte Spiel habe ich mir selbst ausgedacht und "Twelve Card Poker" getauft. 
Dies geschah, da mir zu diesem Punkt noch keine weiteren Forschungsvarianten bekannt waren und ich eine weitere Komplexitätsstufe benötigte. 
Im Verlauf meiner Arbeit habe ich die Varianten Rhode Island Hold'em und Royal Hold'em kennen gelernt und ebenso umgesetzt.
Zusätzlich wurde auch Limit Hold'em implementiert, welches die höchste Komplexität aller in dieser Arbeit Implementierten Varianten aufweist.


\begin{comment}
\chapter{Analyse}

\section{Spielvarianten}
\subsection{Geplante Spielvarianten}
[Was war im Exposé geplant - Liste der Varianten]

\subsection{Implementierte Spielvarianten}
[Was wurde tatsächlich implementiert]

\subsection{Begründung der Auswahl}
[Gründe für die Auswahl dieser Varianten - unterschiedliche Komplexitäten, 
Vergleichbarkeit, etc.]

\section{CFR Solver}
\subsection{Geplante CFR Varianten}
[Was war im Exposé geplant - Vanilla CFR, CFR+, MCCFR, etc.]

\subsection{Implementierte CFR Varianten}
[Was wurde tatsächlich implementiert]

\subsection{Begründung der Auswahl}
[Warum diese Algorithmen - Baseline, Vergleich, Skalierbarkeit, etc.]

\section{Architektur Entscheidungen}
[Game Engine als separate Entität, Erweiterbarkeit, 
Trennung von Algorithmus und Spiel Logik]

\section{Evaluationsmethoden}
\subsection{Geplante Evaluationsmethoden}
[Was war im Exposé geplant]

\subsection{Implementierte Evaluationsmethoden}
[Was wurde tatsächlich implementiert]

\subsection{Begründung der Auswahl}
[Warum diese Methoden - Vergleichbarkeit, Reproduzierbarkeit, etc.]
\end{comment}
